% Chapter 1 -- Starting and Stopping (Part I, Learning; AN-14 task form).
% Sources (Law 1): os/milton/kmain.c (the flagship desktop bring-up: the
% INITECH volume mounted as drive A, TRASH and DESKTOP.DB created when
% missing, FINDER then NOTES then HELLO launched, HELLO foreground; the
% Special-menu power commands finder_power / finder_restart_now /
% finder_shutdown_screen; Exit to InitechDOS finder_exit_dos and the
% transient COMMAND.COM hand-off); os/flair/shell.h (the two stacked menu
% bars); os/flair/menu.c (the menu-bar clock, MenuClock_*);
% os/flair/finder_desktop.c (the volume and Trash icons);
% spec/flair_tenants_demo.h (HELLO 60,60-360,260; NOTES 260,120-560,340);
% spec/dos_banner.txt.  Model: Macintosh User's Guide 1991 Ch. 1 "Turning
% the Macintosh On and Off" (apple-macintosh.md Sec 8.2); Mac OS 8
% Installation task steps (Sec 5.4).
% GATES: test-flair-desktop test-flair-appswitch test-flair-menubar-clock test-flair-restart test-flair-shutdown test-flair-exit-to-dos test-pubs-flair-screens
\Gates{test-flair-desktop, test-flair-appswitch, test-flair-menubar-clock,
  test-flair-restart, test-flair-shutdown, test-flair-exit-to-dos,
  test-pubs-flair-screens}
\chapter{Starting and Stopping}

This chapter shows you how to start InitechOS and FLAIR, what you see when
FLAIR is ready, and the three ways to finish your work: shutting the
computer down, restarting it, and leaving FLAIR for the InitechDOS command
prompt.

\section{Before you begin}
Your computer must be set up and connected as described in
\doctitle{Getting Started with FLAIR}.  The two InitechOS disks must be in
place: the \term{startup disk}, which holds InitechDOS and FLAIR, and the
\term{data disk}, labelled INITECH, which holds your files and the
applications.

\section{Starting FLAIR}
\index{starting FLAIR}\index{FLAIR!starting}
\begin{TaskSteps}
\Step Turn on the display.
\Step Turn on the computer.
\Result{The computer checks its memory and loads InitechOS from the
startup disk.  InitechOS starts FLAIR at once; you do not type a command.
After a few seconds the FLAIR desktop fills the screen.}
\end{TaskSteps}

\FlairScreen[label=fig:desktop]{desktop}{The FLAIR desktop when FLAIR starts}

Figure~\ref{fig:desktop} shows the screen when FLAIR is ready for work.
The parts of it are introduced here and described fully in Chapter~2.

\SideHead{Desktop}The \term{desktop} is the teal working area that fills
the screen.  Everything else lies on it.

\SideHead{Menu bars}Across the top of the screen are two \term{menu bars}.
The upper bar belongs to FLAIR itself: it holds the Apple menu (the apple
at the left), the menu titles \menuname{File}, \menuname{Edit},
\menuname{View} and \menuname{Special}, and the clock.  The lower bar
belongs to the application you are working with, and changes when you
change applications.

\SideHead{Clock}The time of day is shown at the right end of the upper
menu bar, in hours and minutes, with \texttt{am} or \texttt{pm}.  FLAIR
reads the time from the computer's clock and changes the display when the
minute changes.\index{clock}

\SideHead{Disk icon}The icon at the top right, labelled \texttt{INITECH},
stands for the data disk.  Opening it shows the files and folders on the
disk (Chapter~3).

\SideHead{Trash}The \term{Trash} icon at the bottom right is where you put
the files and folders you want to throw away (Chapter~6).

\SideHead{Windows}Two windows are open when FLAIR starts, \texttt{HELLO}
and \texttt{NOTES}.  They belong to two small applications built into
FLAIR, which you can use to practise working with windows (Chapter~3).
\texttt{HELLO} is in front, so the lower menu bar is HELLO's.

\SideHead{Pointer}The arrow is the \term{pointer}.  It moves on the screen
when you move the mouse.  It appears the first time you move the mouse.

\section{Using the mouse}
\index{mouse}\index{pointer}
You give FLAIR most of its instructions with the mouse.  Keep the mouse on
a flat surface beside the keyboard, with its cable pointing away from you.
\begin{TaskSteps}
\Step Move the mouse across the desk.
\Result{The pointer moves in the same direction on the screen.  If you run
out of room on the desk, lift the mouse, put it down where you have room,
and move it again; the pointer does not move while the mouse is in the
air.}
\Step Move the pointer to the \texttt{INITECH} icon.
\Step Press the mouse button and release it.
\Result{This is called \term{clicking}.  The icon is darkened: it is now
\term{selected}.}
\Step Click a bare part of the desktop.
\Result{The icon is no longer selected.}
\end{TaskSteps}
FLAIR uses one mouse button.  If your mouse has two or three buttons, use
the left one; FLAIR ignores the others.\index{mouse button}

\section{Shutting down the computer}
\index{shutting down}\index{Shut Down command}
When you have finished working, shut the computer down before you turn it
off.  Shut Down saves the arrangement of your desktop on the data disk, so
that the windows and icons are where you left them the next time you start
FLAIR, and asks an open application to quit first.
\Operation{FLAIR-SHUTDOWN-SAFE}
\begin{TaskSteps}
\Step Click a bare part of the desktop.
\Result{The lower menu bar changes to the menu bar of the Finder, the part
of FLAIR that looks after the desktop.  Its \menuname{Special} menu holds
the commands that stop the computer.}
\Step Move the pointer to \menuname{Special} in the lower menu bar, and
hold down the mouse button.
\Result{The Special menu drops down.  Its commands appear while you hold
down the button.}
\end{TaskSteps}

\FlairScreen[scale=4, label=fig:special]{specialmenu}{The Special menu}

\begin{TaskSteps}[resume]
\Step Still holding down the button, move the pointer to \menuitem{Shut
Down}, and release the button.
\Result{If an application is open, FLAIR asks it to quit; an application
with unsaved work asks you whether to save it (Chapter~5).  Then the screen
goes black and a message tells you that it is safe to switch off the
computer.}
\end{TaskSteps}

\FlairScreen[label=fig:shutdown]{shutdown}{The message after Shut Down}

\begin{TaskSteps}[resume]
\Step Turn off the computer, then the display.
\Result{If you want to go on working instead, click \texttt{Restart} or
press \key{Enter}: the computer starts again, as if you had just turned it
on.}
\end{TaskSteps}

\begin{MacWarning}
Do not turn off the computer while FLAIR is running.  The arrangement of
your desktop is saved only by Shut Down, Restart and Exit to InitechDOS; a
document you have not saved in its application is lost.
\end{MacWarning}

\section{Restarting the computer}
\index{restarting}\index{Restart command}
Restart does everything Shut Down does and then starts the computer again
at once.
\begin{TaskSteps}
\Step Click a bare part of the desktop.
\Step Choose \menuitem{Restart} from the \menuname{Special} menu.
\Result{FLAIR asks an open application to quit, saves the arrangement of
the desktop and restarts the computer.  FLAIR does not ask you to confirm
Restart: choosing the command starts it.}
\end{TaskSteps}

\section{Leaving FLAIR for InitechDOS}
\index{InitechDOS, leaving FLAIR for}\index{Exit to InitechDOS command}
Some programs, such as the InitechDOS utility programs and the TPS Report
Generator, run only at the InitechDOS command prompt.  To use them, leave
FLAIR.
\Operation{FINDER-EXIT-DOS-ALERT}
\begin{TaskSteps}
\Step Click a bare part of the desktop.
\Step Choose \menuitem{Exit to InitechDOS} from the \menuname{Special}
menu.
\Result{An alert asks you to confirm.}
\end{TaskSteps}

\FlairScreen[label=fig:exitalert]{exitalert}{Confirming Exit to InitechDOS}

\begin{TaskSteps}[resume]
\Step Click \texttt{Exit}, or press \key{Enter}.
\Result{FLAIR asks an open application to quit and saves the arrangement
of the desktop.  Then the desktop disappears and the InitechDOS screen
takes its place.  Because the data disk has no \file{AUTOEXEC.BAT} file,
InitechDOS asks for the date and the time.  To cancel instead, click
\texttt{Cancel} or press \key{Esc}; nothing changes.}
\end{TaskSteps}

\FlairText[fig:exitdate]{exitdate}{InitechDOS after Exit to InitechDOS}

\begin{TaskSteps}[resume]
\Step Press \key{Enter} to keep the date, and \key{Enter} again to keep
the time.
\Result{The command prompt \texttt{A:\textbackslash>} appears.  Drive A is
the data disk; the programs that run at the prompt are in its folder
\file{APPS}.  The \doctitle{InitechDOS User's Guide} tells you how to use
the prompt.}
\end{TaskSteps}
FLAIR does not start again from the command prompt.  To return to FLAIR,
restart the computer: press the computer's reset button, or turn it off
and on again.
